
\section{conclusion}

\begin{frame}
\frametitle{Spin-Polarization Control of Photoelectrons
Using Poincar\'e Fields \\
\small (Phys. Rev. Lett. \textbf{133}, 043201 (2024))}
\framesubtitle{Conclusion}


\small
\begin{textblock*}{400pt}(30pt,60pt)
\vspace{30pt}
\begin{itemize}
\footnotesize
\item Distribution of polarization states 
and angular momentum density of Poincar\'e fields can be used to produce
structured electron helicity distribution
\item This has been demonstrated numerically using an \textit{ab initio}
approach based on the time dependent Dirac equation
\item This \textit{ab initio} approach supersedes what is possible through 
$S$-matrix calculations, which are limited without an analytical solution 
to the Dirac equation
%\begin{itemize}
%\footnotesize
%\item[$\bullet$] Rare quantitative information can be probed about 
%various high intensity effects
%\item[$\bullet$] Importance of long range electronic potential 
%and nonadiabatic effects are better rendered 
%\end{itemize}

\vspace{20pt}

\end{itemize}

\end{textblock*}


%\begin{textblock*}{100pt}(260pt,70pt)
%\includegraphics[scale=.30]{images/medium1.png}
%\end{textblock*}

\end{frame}
